
\subsubsection{Recursive \emph{Proof of Reflection}}

\mk{Maybe recursive PoR stuff should be moved to the PoR section, but that doesn't feel right b/c we don't use it before now (besides dapp-chains but that's trivial).}

\todo{edit start of this section to lead in from previous instead of being burried in Key Concepts}

Part of our security architecture will be the idea of \emph{recursive PoR}.

Say we have three chains: \cL, \cM, and \cR.
They are doing mutual reflection like this: $\xymatrix@C=1em@M=0em{{\cL} \; \ar@{<->}[r] & \; {\cM} \; \ar@{<->}[r] & \;{\cR}}$.
That is: \cL reflects \cM but not \cR; \cM reflects both \cL and \cR; and \cR reflects \cM but not \cL.

With what we already know, it's clear that \cM is as secure as all three combined.
However, that isn't the case for \cL and \cR.
How can we make \cL and \cR as secure as \cM?
Is this even reasonable?

Let's consider chain \cL -- what happens if \cL's local chain-work is, say, $3\times$ that of chain \cM?
The total chain-work that \cL takes into account is $\nicefrac43$ of its local chain-work, so an attacker would need $\nicefrac23$ worth of \cL's local chain-work to control \cL's history.
This is more than 50\%, but it's within reason that an attacker might achieve the required hash-rate.

Case 1: If \cR's local chain-work is large (say, $3\times$ that of \cM), then \cM is almost twice as hard to attack as \cL.\footnote{
    \cM is $\nicefrac74\times$ harder to attack than \cL.
}
This is a problematic situation: it means that an attacker might be able to perform a doublespend against \cL even though they could not perform one against \cM.
Therefore, this case would not be $O(n)$ secure.

Case 2: However, if \cR's local chain-work is small (say, $\nicefrac13$ that of \cM), then \cR is not particularly significant.
Thus, \cM is not substantially harder to attack than \cL, and an attacker could attack both \cL and \cM simultaneously (optionally, \cR too).
But, \cR \emph{is} substantially easier to attack than either \cM or \cL.
The total chain-work that \cR takes into account is only $\nicefrac4{13}$ that of the whole network.
So, this case is \emph{also} not $O(n)$ secure, even though \cL \emph{is} $O(n)$ secure.

These two problem cases are what recursive PoR solves.

Notice that in both cases, \cM is $O(n)$ secure.
Given this, chains \cL and \cR should be able to use the same work that chain \cM does to reach the same level of security.

The first problem is that \cL and \cR cannot use the same \textsc{WeightOf} algorithm as before (although, \cM can).
Particularly, \cL and \cR's \textsc{WeightOf} must count: not only their local chain-work and that contributed reflected \cM blocks, but also some additional weight based on reflections from \emph{other} chains.
Each \cM block might record the total chain-work that it adds to \cM (including PoRs), but \cL and \cR's new \textsc{WeightOf} can't use this value directly.
Firstly, that value includes past local chain-work that reflected \cM, but more importantly: we don't want to force \cL and \cR to verify \emph{every} PoR that contributes to \cM, just the ones that matter.
So how can \cL and \cR nodes know which work to count?

There is a trivial solution to this problem: require \cL and \cR nodes to download the headers of \emph{all} relevant chains, verify their PoRs, and track the origin of those PoRs.
Nodes must track the origin of work contributed via PoR to avoid double-counting work.
An easy example of this is: \cL nodes should not count PoRs where an \cL block is the reflecting block; that work is already directly included in \cL's local chain-work.

% commands used in the following xy chains diag

\providecommand\BL{}
\renewcommand\BL[1]{\BCI{\cL}{i}{#1}}
\newcommand\xBL[1]{\xyB{\BL{#1}}}

\providecommand\BR{}
\renewcommand\BR[1]{\BCI{\cR}{k}{#1}}
\newcommand\xBR[1]{\xyB{\BR{#1}}}

\providecommand\BM{}
\renewcommand\BM[1]{\BCI{\cM}{j}{#1}}
\newcommand\xBM[1]{\xyB{\BM{#1}}}

% ARrow reFlection
\newcommand\arf[1][]{\ar@{.>}[#1]}
\newcommand\arff[1][]{\ar@2{.>}@[Maroon][#1]}

Chain \cL can only count work contributed by \cR \emph{after} a multi-stage PoR can be constructed.
For example, let's consider the (rather orderly) chain-segments in \autoref{fig:simple-recursive-por-chain-diag}.

\begin{figure}
\begin{equation*}
    \xyMChains{
        \xyL{\text{Chain \cL}} & \cdots & \xBL1 \ar[l] & & \xBL2 \ar[ll] \arf[ld] & & \xBL3 \ar[ll] \arf[ld] \\
        \xyL{\text{Chain \cM}} & \cdots & & \xBM1 \ar[ll] \arf[ld] \arf[lu] & & \xBM2 \ar[ll] \arf[lu] \arf[ld] & \\
        \xyL{\text{Chain \cR}} & \cdots & \xBR1 \ar[l] & & \xBR2 \ar[ll] \arf[lu] & & \xBR3 \ar[ll] \arf[lu] \\
        \xyTimeHoriz{rrrrrr} &&&&&&
    }
\end{equation*}
\caption[
    asdf.
]{
    asdf.
}
\label{fig:simple-recursive-por-chain-diag}
\end{figure}

We can see that \BR{2} implicitly images \BL{1} via \BR{2}'s explicit imaging of \BM{1}.
This is because the network already depends on the PoRs that are sufficient to prove this. \BL{2}'s PoR for \BL{1} will prove $\xyInlineRefl{\BL{1}\;&\;\BM{1}\arf[l]}$ (i.e., that \BL{1} was imaged by \BM{1}).
\BM{2}'s PoR (via \cR) for \BM{1} will prove $\xyInlineRefl{\BM{1}\;&\;\BR{2}\arf[l]}$.

Similarly, we know that \BL{3} implicitly images \BR{2}: $\xyInlineRefl{\BR2\;&\;\BM2\arf[l]\;&\;\BL3\arf[l]}$.

Thus, for \BL3, we can construct a full PoR of \BL1 via \BR2 using the proofs for each link between them: $\xyInlineRefl{\BL1 \;&\; \BM{1}\arf[l] \;&\; \BR2\arf[l] \;&\; \BM2\arf[l] \;&\; \BL3\arf[l]}$.

Since \cL nodes already know about all the \cM and \cR blocks (and the relevant PoRs), the full recursive PoR \emph{does not need to be explicitly recorded}.
It is already known to all \cL nodes, since each part of the recursive PoR is available in one of the \cL, \cM, or \cR chains (it is explicitly recorded if a +PoRs UT variant is used, and recalculated by the node otherwise).

In \autoref{sec:counting-work} we discussed the idea that a PoR can only contribute to local blocks that are in the reflecting block's past.
By inspection, we can see that the most recent \cR block (the reflecting block) in \BL{3}'s past is \BR{2}, and the most recent \cL block (the local block) in \BR2's past is \BL1.
Thus, we can see that our method for attributing reflected work -- following the arrows -- works for recursive PoR, too.

There is something important to point out here: \emph{latency}.
Notice that these 4 blocks (in chronological order) are required for this recursive PoR to exist: \BM1, \BR2, \BM2, and \BL3.
This is twice the number of blocks that would normally be required without any recursion (e.g., a valid non-recursive PoR for \BL1 exists after \BM1 and then \BL2 are mined).
Therefore, the delay required for this PoR is twice what we would expect from non-recursive PoR.
In general (if only one PoR path between two given blocks will exist) this delay, measured in block periods, is the length of the full PoR.\footnote{
    For chains with equal block periods (and ignoring \emph{miner resonance}), there is a 50\% chance that chain \cR will produce a block before the next \cL block.

}

Recapping the first problem: \cL and \cR nodes need an expanded \textsc{WeightOf} function with access to \emph{the origin} of implicitly reflected work, and knowledge of \emph{where} that work has already been counted.
Nodes can do this relatively efficiently, for all relevant chains, by downloading the headers and verifying the PoRs.
This method does not require recursive PoRs to be explicitly recorded in full since the components of these PoRs are already available and have already been individually verified.
Our existing rules about how to attribute the weight of PoRs (follow the arrows) still works.

Our solution to the first problem is promising, but we have used an implicit assumption: there is only one viable path for PoRs between \cL and \cR.
What if there are \emph{multiple paths} between \cL and \cR?
This is the second problem.

\subsubsubsection{Recursive PoR with Multiple Paths}

\providecommand\BMone{}
\renewcommand\BMone[1]{\BCI{\cMone}{g}{#1}}
\newcommand\xBMone[1]{\xyB{\BMone{#1}}}

\providecommand\BMtwo{}
\renewcommand\BMtwo[1]{\BCI{\cMtwo}{h}{#1}}
\newcommand\xBMtwo[1]{\xyB{\BMtwo{#1}}}

% ARrow reFlection Bent Down
\newcommand\arfbd[1][]{\ar@{.>}@/_9pt/[#1]}
\newcommand\arfbu[1][]{\ar@{.>}@/^9pt/[#1]}

Say that we now have four chains: \cL, \cMone, \cMtwo, and \cR doing mutual PoR like so:
\begin{equation*}
    \xymatrix@C=2em@R=-.15em{
        & {\cMone}
        \ar@{<->}[dd]
        & \\
        {\cL} \ar@{<->}[ur] & & {\cR} \ar@{<->}[ul] \\
        & {\cMtwo} \ar@{<->}[ul] \ar@{<->}[ur] &
    }
\end{equation*}
Both \cMone and \cMtwo are already $O(n)$ secure.
Assuming that \cL and \cR nodes download and verify all PoRs, how can \cL and \cR count each other's work?

First, let's consider when \cL can \emph{definitely} count \emph{all} the work contributed by a \cR block.

\begin{figure}
    \begin{equation*}
        \xyMChains{
            \xyL{\text{Chain \cL}} & \cdots & \xBL1 \ar[l] & & \xBL2 \ar[ll] \arf[ld] \arfbu[ldd] & & \xBL3 \ar[ll] \arf[ld] \arfbu[ldd] \\
            \xyL{\text{Chain \cMone}} & \cdots & & \xBMone1 \ar[ll] \arfbd[ldd] \arf[lu] & & \xBMone2 \ar[ll] \arf[lu] \arfbd[ldd] & \\
            \xyL{\text{Chain \cMtwo}} & \cdots & & \xBMtwo1 \ar[ll] \arf[ld] \arfbu[luu] & & \xBMtwo2 \ar[ll] \arfbu[luu] \arf[ld] & \\
            \xyL{\text{Chain \cR}} & \cdots & \xBR1 \ar[l] & & \xBR2 \ar[ll] \arf[lu] \arfbd[luu] & & \xBR3 \ar[ll] \arf[lu] \arfbd[luu] \\
            \xyTimeHoriz{rrrrrr} &&&&&&
        }
    \end{equation*}
    \caption[
        asdf.
    ]{
        asdf.
        Note that reflections between \cMone and \cMtwo are not shown.
    }
    \label{fig:multi-path-recursive-por-chain-diag}
\end{figure}

As before, we'll look at valid PoRs for \BL3 that prove \BL1 was reflected in \BR2.
The exact reflections that exist will be analogous to our prior example, too.
See \autoref{fig:multi-path-recursive-por-chain-diag} for the chain-diagram.
% For each $x$ in $\{1,2,3\}$, we'll assume that \BL{x} and \BR{x} occur simultaneously, and both image \BMone{x-1} and \BMtwo{x-1}.
% After \BL{x}, but before \BL{x+1}, we'll assume that \BMone{x} and \BMtwo{x} are mined simultaneously and both blocks image \BL{x} and \BR{x}.
Notice that there are 4 possible paths from \BL{3} to \BL1:
% \\[-.66em]
\begin{equation*}
    \xymatrix@C=1em@R=-.5em{
        & & \BMone1 \arf[ld] & & \BMone2 \arf[ld] & \\
        \BL0 & \BL1 & & \BR2 \arf[lu] \arf[ld] & & \BL3 \arf[lu] \arf[ld] & \BL4 \\
        & & \BMtwo1 \arf[lu] & & \BMtwo2 \arf[lu] & \\
    }
\end{equation*}
% \\[-.66em]
Here is what we can say about this case: if all 4 of these PoR paths exist, then \BR2 \emph{definitely} reflects \BL1, and its weight can be included in \BL3's total chain-weight.

But, what if only some of those paths are available?
Consider this alternative case:
% hmm, unsure about the colors / focused arrows
\begin{equation*}
    \xymatrix@C=1em@R=-.15em{
        & \BMone1 \arff[ld] & & & \BMone2 \arf[llld] \arff[ld] & \\
        \BL0 & \BL1 & & \BR2 \arff[llu] \arff[ld] & & \BL3 \arff[lu] \arf[llld] & \BL4 \arf[ld] \\
        & & \BMtwo1 \arff[lu] & & & \BMtwo2 \arf[llu] \\
    }
\end{equation*}
In this case, \cL is implicitly reflected in \BR2, and \BL3 has two possible recursive PoRs of this (indicated via the double-lined arrows): \BL{1} via \BMtwo1; and \BL{0} via \BMone1.
Are both valid?
Is either preferable?
If \BL3 includes the full weight of \BR2 in its chain-weight, then what happens when \BL4 is mined and the other two PoR paths via \BMtwo2 become available?

We \emph{can} be sure that \BR2's weight can be fully included in \cL's chain-weight \emph{after} \BL4 is mined (as all PoR paths now exist).
However, it's not so clear what we should do when not all PoR paths yet exist.
Do we even need to solve this problem, now, though?

\begin{itemize}
    \item no -- it's enough to know that, when all PoR paths exist, the full weight is okay
    \item this is b/c we'll simplify tiles: every ${B_f}^{-1}$ seconds all chains in each tile mine 1 new block and provide reflections of 1 new block in their local tile and adjacent tiles, and vice versa.
    Thus all PoR paths are available all at once.
    \item we just need to be confident that an answer \emph{does} exist
    \item also, simplification case is much easier to reason about and solve for
    \item full solution left as a future research
    \item some ideas or questions for that future research:
    \begin{itemize}
        \item is a single PoR path enough to count the full weight of the reflecting block?
        \item if so -- tiling is very fast and secure. if not, tiling is pretty fast and still secure.
        \item what happens if some paths \emph{never} end up existing?
        \item do we divide contributed work by the number of paths?
        \item do paths have to go through the same chain? (i.e., $\cL \to \cMone \to \cR \to \cMone \to \cL$)
        \item can paths go through different chains? (i.e., $\cL \to \cMone \to \cR \to \cMtwo \to \cL$)
        \item should each PoR be valid if 1st and last entries are removed?
        if so, then it passes its weight to like that chain-segment (of the intermediate chain for which the middle bits of the recursive PoR are valid).
        \item but PoRs via different chains should still be valid b/c it shows order (thus we know that whatever's being reflected is in the history of every intermediate block)
        \item do we ever really need to consider this partial paths stuff?
        like a block (\BL3) that enables some PoR paths via a new block (\BR2) will probs also enable other missing PoR paths.
        Example: \BL4 enables \BR2 via \BMtwo2, but could also mb enable \BR3 via \BMtwo2, too.
        We should expect all paths to be available within a few block periods (e.g., $\sim 60$ s).
        \item WRT paths like $\cL \to \cMtwo \to \cMone \to \cR \to \cMone \to \cMtwo \to \cL$ (i.e., it's longer than the shortest path possible) -- we shouldn't need to count these.
        we should always be able to go directly b/c an implicit reflection in a chain's history that can also be explicit, should be explicit.
        so count these as 0 weight.
        \item idea: only the paths back matter. making more paths between the \cR block and most recent \cL block isn't so hard, but making the \cR block is (if the \cR block is relatively heavy, that is).
        So: divide \cR block weight by the number of \emph{paths back} and then only count unique shortest paths from $\BR{x} \to \BL{y}$.
        By that point, shouldn't \emph{all} those paths be available, anyway?
        \item maybe: path availability is only an issue between the reflecting \cR block and the current \cL block, but not between the \cR block and the reflected \cL block?
        (This sounds like it logically has to be true -- how could any new paths back be introduced!?)
        \item idea: the weights of all blocks in a recursive PoR path \emph{all} have the same \cL in their history, so their weights can be safely added to \cL.
        \item idea: the principle \emph{don't double-count work} means that when we have multiple paths back, that block's weight is attributable to the \emph{most recent} reflected \cL block.
        \item idea: the principle that \emph{a PoR is worth whatever is required to attack it} means that only 1 PoR path to the \cR block is required -- once that happens all other paths are redundant (though useful for SPV).
        \item idea: so the total weight contributed by a set of PoRs is the sum of all unique blocks in that set of PoRs -- multiple paths doesn't help.
        \item so simplex-chains in a deep tile get the security benefit of root tile chains as soon as a single path is available from an L block to any R block back to any L block.
        \item this should happen MUCH faster than the predicted delay between singularly reflected chains -- $v\cdot{N_1}^{-1}\times$ the duration, I think.
    \end{itemize}
    \item do miners need to be incentivized to reflect child chains?
\end{itemize}

There are some future areas of research that may improve this process, too.
For example: along the lines of NIPoPoWs\footnote{
    Non-Interactive Proofs of Proof-of-Work (NIPoPoWs), see \citeNIPoPoWsFull.
}, are \emph{Non-Interactive Proofs of Proof-of-Work Reflections} (NIPoPoWRs) possible?

\todo{por sim: are we double-counting reflections sometimes? should we divide the weight of a block by the number of local blocks that in reflects? what about blocks that have no explicit reflections? need to check sim logic.}

\todo{continue}

\subsubsubsection{Principles of PoR}

\begin{itemize}
    \item work can be converted to common units and thus summed
    \item work should not be double counted
    \item work is attributable to the most recent L block reflected by an R block (implies: pow's weight is transitive to all things in that PoWs history)
    \item work is work (law of identity; it means that it doesn't matter where/how work happens, just that it does and that its history includes the thing we want to prove)
\end{itemize}
